\section{Resultados}

\subsection{Visualización conjunta de las estrategias}

La Figura~\ref{fig:scatter} presenta el frente de Pareto aproximado por
NSGA-II (círculos azules) junto con las cuatro soluciones seleccionadas por
las distintas estrategias. Esta gráfica responde al \textbf{primer entregable}
del enunciado: la comparación espacial de las decisiones con la leyenda
visible. El frente puede explorarse de forma interactiva (moviendo los pesos y
el criterio lexicográfico) en el dashboard publicado en
\href{https://multicriterio.streamlit.app}{multicriterio.streamlit.app}.

\begin{figure}[h]
  \centering
  \includegraphics[width=\linewidth]{scatter_estrategias.pdf}
  \caption{Comparativa de estrategias de decisión multicriterio. El frente de
  Pareto (NSGA-II) define la curva de soluciones de compromiso. Las dos
  soluciones \emph{ponderadas} (a priori y a posteriori) caen prácticamente
  superpuestas en el ``codo'' del frente, mientras que ambas soluciones
  \emph{lexicográficas} coinciden en el extremo de menor tiempo
  $(f_1, f_2) \approx (\num{200}, \num{0})$.}
  \label{fig:scatter}
\end{figure}

\subsection{Soluciones seleccionadas}

La Tabla~\ref{tab:soluciones} reporta las cuatro soluciones. Todos los valores
se leen directamente de \texttt{results/\allowbreak selected\_solutions.json}.

\begin{table}[h]
  \centering
  \caption{Soluciones obtenidas por cada estrategia de decisión.}
  \label{tab:soluciones}
  \begin{tabular}{llSS}
    \toprule
    \textbf{Enfoque} & \textbf{Criterio} & {\textbf{$f_1$ (Costo)}} & {\textbf{$f_2$ (Tiempo)}} \\
    \midrule
    A Posteriori & Suma Ponderada  & \PostPondCosto    & \PostPondTiempo \\
    A Posteriori & Lexicográfica   & \PostLexCosto     & \PostLexTiempo \\
    A Priori     & GA Ponderado    & \AprioriPondCosto & \AprioriPondTiempo \\
    A Priori     & GA Lexicográfico& \AprioriLexCosto  & \AprioriLexTiempo \\
    \bottomrule
  \end{tabular}
\end{table}

\subsection{Calidad del frente de Pareto}

El frente aproximado por NSGA-II consta de \MetNSol{} soluciones no dominadas,
con los siguientes indicadores (de \texttt{results/metrics.json}):

\begin{itemize}[noitemsep]
  \item \textbf{Hipervolumen} (ref.\ $(\num{200},\num{50})$):
        \num{\MetHipervolumen}.
  \item \textbf{Spacing} de Schott: \num{\MetSpacing} (valor bajo $\Rightarrow$
        distribución uniforme).
  \item \textbf{Extremos del frente:} $f_1 \in [\num{\MetMinFuno},\ \num{\MetMaxFuno}]$,
        $f_2 \in [\num{\MetMinFdos},\ \num{\MetMaxFdos}]$.
\end{itemize}

Los extremos coinciden exactamente con la predicción analítica: el frente va
de $(\num{0}, \num{50})$, correspondiente a $\mathbf{x}=(0,0)$ (mínimo costo),
hasta $(\num{200}, \num{0})$, correspondiente a $\mathbf{x}=(5,5)$ (mínimo
tiempo). Esta concordancia valida numéricamente la implementación.
